\section{Related work}\label{sec:related}

\paragraph{Failure attribution in multi-agent LLM systems.}
A growing line of work asks which agent, and at which step, caused a task to fail, given the log of one multi-agent run. Who\&When~\citep{zhang2025which} introduced a benchmark and studied LLM judges for the task; AgenTracer~\citep{zhang2025agentracer} trains a model to attribute failures; hierarchical and consensus-based approaches improve step-level accuracy~\citep{banerjee2025where}; and FALAT~\citep{rafi2026falat} frames attribution as a dependency-guided search over decisions, tool outputs and messages. These methods assume that the trajectory is recorded in full and ask about one failing run. Our setting differs in two ways: we ask how an item \emph{spread} among many agents through shared state, and the thing that links agents, a read, is exactly what is missing from the log.

\paragraph{Provenance and information-flow tracking.}
Data provenance~\citep{buneman2001why} and information-flow control~\citep{denning1976lattice,enck2010taintdroid} track how data moves through a system that is instrumented to follow it. For LLM agents, recent work tracks information flow through contexts, tool calls and memory~\citep{cai2026ghost}, and a survey organises evidence tracing and execution provenance for agents~\citep{wang2026traces}. Instrumentation is the right answer when the operator controls the data path. Our contribution sits on both sides of that line: an audit that says how much can be recovered when nothing was instrumented, calibrated against the coincidences that shared model habits create, and a gateway design whose tokens make the instrumentation survive being copied by an agent that is free to rewrite the text.

\paragraph{Planted markers.}
Fictitious map entries~\citep{fictitiousentry}, canary tokens~\citep{canarytokens} and traitor tracing~\citep{chor1994tracing} all plant a recipient-specific mark and read it back from a leak. For agents, canary-style protocols have been used to trace an agent back to its owner~\citep{chocron2026owns}, and keyed signals in a model's token distribution to audit which agent produced a text~\citep{nian2026final}. These identify who \emph{produced} an output. We use the same family of ideas to identify which \emph{served copy} an agent was exposed to, which is the quantity that matters when a mistake spreads through shared artifacts, and we study empirically when such a mark survives an LLM's tendency to copy links verbatim or to re-derive them.

\paragraph{Contagion and source inference.}
Locating the source of a rumour or a cascade from partial observations is a classical problem~\citep{shah2011rumors,gomezrodriguez2010inferring}. Those methods assume that the spread is observed through a graph or through activation times. In a swarm that shares artifacts, what is observed are edits, the graph is the unobserved reads, and the ceiling on what is recoverable is a function of what is logged (Section~\ref{sec:theory}).

\paragraph{Mistakes and attacks that spread among agents.}
Infectious jailbreaks and prompt infection show that a single poisoned input can propagate through a population of agents~\citep{gu2024agentsmith,lee2024prompt}, which is why post-incident tracing matters for safety as well as for debugging.

\paragraph{Shared habits.}
Models trained on similar data make similar choices, a form of algorithmic monoculture~\citep{kleinberg2021algorithmic}. For tracing, monoculture is the reason a repeated value is weak evidence of copying, and Proposition~\ref{prop:calib} bounds the cost of ignoring it.

\paragraph{Position.}
To our knowledge no earlier study measures, on real multi-agent logs and in a controlled lab, how much copying between agents is traceable when reads are not logged, calibrates copy calls against shared habits, and evaluates a tracing token across model families with a hidden read log as ground truth. We have searched but not exhaustively; related work may exist that we have missed.
